\clearpage
\section{Historically Superseded Postconciliar Forms (1967--2011)}
Postconciliar liturgical law developed across three typical editions of the Roman Missal (1970, 1975, and 2002/2008) together with papal motu proprios and dicasterial instructions. The three chronological entries below record provisions and practices that occurred during earlier postconciliar decades and have since been abolished, clarified, or explicitly prohibited. They are documented here strictly as historical stages and never as current liturgical law.

\label{pc-obj-hist-subdeacon-tunicle-maniple}
\PCDHistHead{Subdeacon's Tunicle and Maniple}{Tunicella et manipulus}{1967--1972 --- Tres abhinc annos (1967) no. 25; Ministeria quaedam (1972); 1975 Missale Romanum; CE (1984) no. 56}{N41}
Between 1967 and 1975, two vestments inherited from the preconciliar sanctuary passed out of postconciliar Roman use in stages. First, the Sacred Congregation of Rites Instruction Tres abhinc annos (4 May 1967, effective 29 June 1967), no. 25, declared that the maniple---the ornamental strip of liturgical cloth worn on the left forearm of the priest, deacon, and subdeacon---was no longer required; it ceased to be prescribed and was omitted from the 1970 General Instruction of the Roman Missal rather than being formally banned by decree. Second, the first typical edition of the postconciliar Missale Romanum (26 March 1970, Prot. 166/70) briefly retained the liturgical office of the subdeacon (1970 IGMR nos. 142--152) and his sleeved outer vestment, the tunicle (no. 301). Pope Paul VI's Motu Proprio Ministeria quaedam (15 August 1972, AAS 64:529--534; effective 1 January 1973) then suppressed first tonsure, the four minor orders, and the major order of subdeacon in the Latin Church, establishing the lay instituted ministries of Lector and Acolyte. The second typical edition of the Missale Romanum (27 March 1975) excised the subdeacon and tunicle from the General Instruction, and the Caeremoniale Episcoporum (14 September 1984) no. 56 retained only a single light pontifical dalmatic beneath the bishop's chasuble on solemnities, omitting the separate pontifical tunicle.

\PCDField{Current discipline (2008/2011)}{The 2002/2008 Latin and 2011 United States Roman Missal prescribe neither the subdiaconate nor the tunicle and maniple (GIRM nos. 119, 335--341). Duties formerly assigned to the subdeacon belong to instituted lectors and acolytes vested in the alb (GIRM no. 336; CIC can. 230, par. 1, as amended by Spiritus Domini, 2021), and on major pontifical occasions the bishop wears only a single light dalmatic beneath the chasuble (Caeremoniale Episcoporum 1984 no. 56).}
\PCDField{1962 comparison}{\PCDTag{superseded} In the 1962 Roman Missal, the maniple was required on the left forearm of the priest, deacon, and subdeacon during Mass, and the tunicle served both as the proper outer vestment of the subdeacon at Solemn Mass and as the inner layered garment beneath the bishop's dalmatic and chasuble. Both vestments passed out of postconciliar law between 1967 and 1975 (and from episcopal layering in 1984), remaining in liturgical use only where the 1962 Missal is celebrated.}
\PCDEntryRule

\label{pc-obj-hist-fragile-vessels-pouring}
\PCDHistHead{Fragile Vessels, Flagon Pouring, and EMHC Purification}{Vasa fragilia et infusio}{1970--2006/2011 --- GIRM (2002/2008) nos. 279, 328--330; Redemptionis Sacramentum (2004) nos. 106, 117, 119; CDWDS letter (Oct 2006); USCCB Norms (2011) no. 53}{N42}
The 1970 and 1975 editions of the General Instruction of the Roman Missal (nos. 290--291) omitted precious metal as an explicit primary norm for sacred vessels, directing only that vessels be fashioned from solid materials esteemed as noble in each region (with preference for materials not easily broken or rendered unusable) and that the chalice bowl be nonabsorbent. Over the next three decades, some parishes adopted glass, stoneware, ceramic, or clay chalices and bowls, and consecrated wine in large flagons or pitchers from which the Precious Blood was poured into Communion cups after consecration. The third typical edition (2002/2008 GIRM nos. 328--330) restored precious metal---gilded inside whenever the metal rusts or is less precious than gold---as the universal primary standard, while the United States adaptation (US GIRM no. 329) narrowed secondary solid noble materials such as hard woods to patens, ciboria, and pyxes that do not easily break, deteriorate, or absorb liquids. On 25 March 2004, the Congregation for Divine Worship and the Discipline of the Sacraments issued Redemptionis Sacramentum, which strictly forbade post-consecration pouring of the Precious Blood and barred flagons or bowls from holding the Blood of Christ (no. 106), while explicitly reprobating glass, earthenware, clay, and fragile or rusting vessels (no. 117). Finally, a temporary three-year United States indult (2002--2005) allowing extraordinary ministers of Holy Communion to purify sacred vessels expired without renewal per the Congregation's letter of October 2006, and the 2011 USCCB Norms (no. 53) strictly reserve purification to the priest, deacon, or instituted acolyte.

\PCDField{Current discipline (2008/2011)}{Sacred vessels are made primarily of precious metal gilded inside where necessary (GIRM no. 328), with nonabsorbent bowls for all chalices (no. 330) and the United States adaptation limiting solid noble alternatives to hard, non-breaking, non-absorbent patens, ciboria, and pyxes (US GIRM no. 329). Glass, earthenware, and clay vessels and post-consecration pouring from flagons are strictly prohibited (Redemptionis Sacramentum nos. 106, 117). Wine brought in a flagon is poured into chalices before the Eucharistic Prayer, and vessel purification belongs exclusively to the priest, deacon, or instituted acolyte (GIRM no. 279; USCCB Norms no. 53).}
\PCDField{1962 comparison}{\PCDTag{no established counterpart} Fragile or porous sacred vessels, post-consecration pouring from flagons, and lay vessel purification had no counterpart in the 1962 Roman Missal, which mandated gold or interior-gilded silver chalices and patens, distributed Communion to the laity under bread alone, and restricted ablutions exclusively to the priest celebrant. The 2002/2008 Missal, Redemptionis Sacramentum (2004), and the 2011 United States norms restore precious metal as the universal primary norm, bar fragile materials and post-consecration pouring, and restrict purification to the priest, deacon, or instituted acolyte.}
\PCDEntryRule

\label{pc-obj-hist-1970-sanctuary-clarifications}
\PCDHistHead{Sanctuary, Cross, Tabernacle, Altar Card, and Veil Clarifications (1970/1975)}{Crux, tabernaculum, tabellae et velum}{1970--2002/2008 --- CDWDS Response (25 Sept 2000, Prot. N. 2036/00/L); GIRM (2002/2008) nos. 118c, 299, 306, 308, 314--315}{N43}
The third typical edition of the Roman Missal (2002/2008) together with a formal Roman clarification in 2000 resolved four sanctuary ambiguities inherited from the 1970 and 1975 editions. First, whereas 1975 GIRM no. 270 spoke only of a cross visible to the assembly without specifying a corpus---producing widespread bare crosses---2002/2008 GIRM no. 308 explicitly mandates a cross bearing the figure of Christ crucified (cum figura Christi crucifixi) on or near the altar, remaining there appropriately even outside Mass. Second, when the freestanding altar wording of GIRM no. 299 raised questions about historic churches, the Congregation for Divine Worship and the Discipline of the Sacraments issued its Response of 25 September 2000 (Prot. N. 2036/00/L; Notitiae 36:171--173), clarifying that expedit expresses desirability rather than an absolute ban on historic wall-attached altars or celebration ad orientem. Third, 2002/2008 GIRM nos. 314--315 refined 1975 GIRM no. 276 by placing two tabernacle locations on equal footing under the judgment of the diocesan bishop: either in the sanctuary apart from the altar of celebration (including on a former altar no longer used for Mass) or in an organically connected, readily noticeable adoration chapel. Fourth, the three 1962 altar cards (tabellae secretarum) were omitted from the 1970 Missal and remain excluded from the mensa under GIRM no. 306, whereas 2002/2008 GIRM no. 118c restored explicit praise of covering the prepared chalice with a veil (color of the day or white), leaving the burse unnamed as a customary accessory.

\PCDField{Current discipline (2008/2011)}{The altar or processional cross must display the figure of Christ crucified (GIRM no. 308); freestanding altars are recommended where possible without banning historic wall altars or celebration ad orientem (GIRM no. 299; CDWDS 25 Sept 2000); the diocesan bishop judges between placing the tabernacle in the sanctuary apart from the altar of celebration or in an organically connected adoration chapel (GIRM nos. 314--315); altar cards are excluded from the mensa (GIRM no. 306); and covering the prepared chalice with a veil of the color of the day or white is explicitly praised, with the burse remaining an unnamed customary accessory (GIRM no. 118c).}
\PCDField{1962 comparison}{\PCDTag{materially changed} The 1962 Roman Missal required a central altar crucifix, ordinarily placed the tabernacle in the middle of the high altar, prescribed three altar cards during Mass, and mandated a silk chalice veil with matching burse. Postconciliar clarifications in 2000--2008 reaffirm the crucifix, permit freestanding or historic wall altars, place sanctuary and chapel tabernacle sites on equal footing under the diocesan bishop's judgment, supersede the three altar cards on the mensa, and commend the chalice veil as a praiseworthy optional practice while leaving the burse as an unnamed custom.}
\PCDEntryRule
